\centerline{\bf \Huge Комбо разнобой}

\begin{flushright}
        \scriptsize{Эпизод 10. Ныряльщик в магме}
\end{flushright}

\problem{1}
В Национальной Баскетбольной Ассоциации 30 команд, каждая из которых проводит за год 82 матча с другими командами в регулярном чемпионате. Сможет ли руководство Ассоциации разделить команды (не обязательно поровну) на Восточную и Западную конференции и составить расписание игр так, чтобы матчи между командами из разных конференций составляли ровно половину от общего числа матчей?

\problem{2}
На встречу ЭлМШ пришли несколько семей (мама, папа и дети), в каждой из которых было от 1 до 10 детей. Расул Владимирович выбирал одного ребенка, одну маму и одного папу из трёх разных семей и давал им задачу на арктангенсы. Оказалось, что у него было ровно 3630 способов выбрать нужную тройку людей. Сколько всего детей могло быть на встрече?

\problem{3}
В стране есть $n>1$ городов, некоторые пары городов соединены двусторонними беспосадочными авиарейсами. При этом между любыми двумя городами существует единственный авиамаршрут (возможно, с пересадками). Мэр каждого города $X$ подсчитал количество таких нумераций всех городов числами от 1 до $n$, что на любом авиамаршруте, начинающемся в $X$, номера городов идут в порядке возрастания. Все мэры, кроме одного, заметили, что их результаты подсчётов делятся на 2016. Докажите, что и у оставшегося мэра результат также делится на 2016.

\problem{4}
В некоторых клетках квадрата $200 \times 200$ стоит по одной фишке - красной или синей; остальные клетки пусты. Одна фишка видит другую, если они находятся в одной строке или одном столбце. Известно, что каждая фишка видит ровно пять фишек другого цвета (и, возможно, некоторое количество фишек своего цвета). Найдите наибольшее возможное количество фишек, стоящих в клетках.

\problem{5}
Даны $n$ монет попарно различных масс и $n$ чашечных весов, $n>2$. При каждом взвешивании разрешается выбрать какие-то одни весы, положить на их чаши по одной монете, посмотреть на показания весов и затем снять монеты обратно. Какие-то одни из весов (неизвестно, какие) испорчены и могут выдавать случайным образом как правильный, так и неправильный результат. За какое наименьшее количество взвешиваний можно заведомо найти самую тяжёлую монету?


